% TOP_PROBLEM: {"id": 109, "title": "Large Sieve and Quadratic Sets", "category_id": 1, "category": {"id": 1, "name": "number_theory", "display_name": "Number Theory", "description": "Properties of integers, prime numbers, Diophantine equations.", "slug": "number-theory", "order_index": 1, "created_at": "2026-07-31T15:26:25.670Z"}, "status": "open"}
% FIRST_POSTED: null
% ENTRY_KIND: statement_only
% Imported from: batch02.tex; UnsolvedMath ID 109
\documentclass[11pt]{article}
\usepackage[margin=25mm]{geometry}
\usepackage{fontspec}
\setmainfont{Latin Modern Roman}
\usepackage{amsmath,amssymb,mathtools}
\usepackage{unicode-math}
\setmathfont{Latin Modern Math}
\usepackage{xurl,hyperref}
\hypersetup{colorlinks=true,urlcolor=blue}
\setcounter{secnumdepth}{0}
\setlength{\parindent}{0pt}
\setlength{\parskip}{5pt}
\setlength{\emergencystretch}{3em}
\newcommand{\sref}[2]{\hyperlink{source:#1:#2}{[S#2]}}
\newcommand{\cataloguescope}[1]{\paragraph{Recorded target.}#1}
\begin{document}
% UnsolvedMath ID 109; source catalogue code GREEN-021
\setcounter{section}{108}
\section{Large Sieve and Quadratic Sets}\label{109}
{\small\sffamily UnsolvedMath ID 109 · Number Theory}
\subsection{Definitions and mathematical statement}

\paragraph{Shared notation.}$\mathbb N=\{1,2,\ldots\}$, and $\mathbb Z,\mathbb Q,\mathbb R,\mathbb C$ denote the integers, rationals, reals and complex numbers. Any other conventions are specified locally.


For $A\subset\mathbb N$ and a prime $p$, write $A\bmod p=\{a+p\mathbb Z:a\in A\}\subset\mathbb Z/p\mathbb Z$. A rational quadratic map is $\phi(t)=at^2+bt+c$ with $a,b,c\in\mathbb Q$ and $a\ne0$; its integer-argument image is $\phi(\mathbb Z)=\{\phi(n):n\in\mathbb Z\}$. Throughout this statement $\log$ is the natural logarithm. The notation $u(X)\ll v(X)$ means that $u(X)\le C v(X)$ for all sufficiently large $X$, for a constant $C$ independent of $X$.
\paragraph{Statement recorded in the supplied source.}
Suppose that $A\subset\mathbb N$ satisfies
\[
|A\bmod p|\le\frac{p+1}{2}\quad\text{for every sufficiently large prime }p.
\]
Must at least one of the following alternatives hold?
\begin{itemize}
\item There exists $C>0$ such that $|A\cap\{1,\ldots,\lfloor X\rfloor\}|\le C X^{1/2}/(\log X)^{100}$ for all sufficiently large $X$.
\item There is a rational quadratic map $\phi$ such that $A\subset\phi(\mathbb Z)$.
\end{itemize}
The second alternative requires containment of the whole set in a single quadratic image, as stated in Problem 47. The relevant quadratic cardinality scale is $X^{1/2}$ for integers up to $X$. \sref{109}{1}\sref{109}{2}
\subsection{Short English statement}
If a set of positive integers occupies at most half the residue classes modulo every sufficiently large prime, must it either be substantially smaller than square-root size or lie entirely among the values of one rational quadratic polynomial at integers?
\subsection{Sources}
\begin{description}
\item[\hypertarget{source:109:1}{S1}] ULAM.AI and the UnsolvedMath Contributors, \emph{UnsolvedMath: A Curated Collection of Open Mathematics Problems}, record 109 (GREEN-021). CC BY 4.0. \url{https://huggingface.co/datasets/ulamai/UnsolvedMath}.
\item[\hypertarget{source:109:2}{S2}] Ben Green, \emph{100 Open Problems}, author manuscript with December 2025 updates, Problem 47 and its comments, p. 23. \url{https://people.maths.ox.ac.uk/greenbj/papers/open-problems.pdf}.
\end{description}
\label{end:109}
\end{document}
